\section*{Chapter \ref{ch:b2:magma-emplacement} --- Magmas and Their Emplacement}

\begin{solution}{exo:b2:magma-emplacement:1}
$T = \qty{1073}{K}$. For $w = 1$, $\ln w = 0$: $\log_{10}\eta = -3.545 +
9601/(1073 - 195.7) = -3.545 + 10.94 = 7.40$. For $w = 4$, $\ln 4 = 1.386$:
$A = -3.545 + 1.155 = -2.390$, $B = 9601 - 3283 = 6318$, $C = 195.7 + 44.7 =
240.4$, and $\log_{10}\eta = -2.390 + 6318/832.6 = -2.390 + 7.588 = 5.20$. The
viscosity falls by $10^{7.40 - 5.20} = 10^{2.20} \approx 160$.
\end{solution}

\begin{solution}{exo:b2:magma-emplacement:2}
Rhyolite: $p_s = (6/0.44)^2 = \qty{186}{MPa}$, at a depth of $\num{186e6}/(2600\times9.81)
= \qty{7.3}{km}$. Basalt: $p_s = (0.5/0.34)^2 = \qty{2.16}{MPa}$, at
$\num{2.16e6}/(2600\times9.81) = \qty{85}{m}$. The rhyolite starts to form bubbles
deep in the crust, the basalt only near the surface (if it carries no CO$_2$).
\end{solution}

\begin{solution}{exo:b2:magma-emplacement:3}
$(1150 + 50)/2 = \qty{600}{\celsius}$ and $(1150 + 400)/2 = \qty{775}{\celsius}$. The
\omterm{def:b2:magma-emplacement:bodies}{dyke} in the cold rock is chilled more strongly at its contact, and has the wider
\omterm{def:b2:magma-emplacement:aureole}{chilled margin}.
\end{solution}

\begin{solution}{exo:b2:magma-emplacement:4}
(a) A \omterm{def:b2:magma-emplacement:bodies}{sill}, parallel to the bedding. (b) A \omterm{def:b2:magma-emplacement:bodies}{dyke}, cutting across the bedding.
(c) A \omterm{def:b2:magma-emplacement:bodies}{batholith}, more than \qty{100}{km^2} and assembled from many intrusions.
(d) A \omterm{def:b2:magma-emplacement:bodies}{laccolith}, which has domed its roof. (e) A \omterm{def:b2:magma-emplacement:bodies}{dyke}: granite has no layering,
so any sheet in it is a \omterm{def:b2:magma-emplacement:bodies}{dyke}, whatever its orientation.
\end{solution}

\begin{solution}{exo:b2:magma-emplacement:5}
$(1 - 1.35\phi)^{-2.5}$: for $\phi = 0.05$, $0.9325^{-2.5} = 1.19$, so
\qty{119}{Pa.s}, against $1 + 2.5\times0.05 = 1.125$, \qty{113}{Pa.s}, with Einstein's
formula (close, as it should be in the dilute limit). For $0.20$: $0.73^{-2.5} =
2.20$, \qty{220}{Pa.s}; for $0.40$: $0.46^{-2.5} = 6.97$, \qty{697}{Pa.s}; for $0.60$:
$0.19^{-2.5} = 63.6$, about \qty{6360}{Pa.s}. Jamming at $\phi = 1/1.35 = 0.74$.
\end{solution}

\begin{solution}{exo:b2:magma-emplacement:6}
$G = 300\times9.81 = \qty{2943}{Pa/m}$. (a) $\bar u = 2943\times1^2/(12\times100) =
\qty{2.45}{m/s}$; $q = \qty{2.45}{m^2/s}$; $Re = 2700\times2.45\times1/100 = 66$, laminar;
$t = 10000/2.45 = \qty{4080}{s}$, about \qty{1.1}{h}. (b) $\bar u =
2943/(12\times\num{e6}) = \qty{2.45e-4}{m/s}$, $t = \qty{4.08e7}{s} = 472$ days; with
$w = \qty{10}{m}$, $\bar u$ is a hundred times larger, \qty{0.0245}{m/s}, and $t =
\qty{4.08e5}{s} = 4.7$ days. (c) The \qty{1}{m} \omterm{def:b2:magma-emplacement:bodies}{dyke} cools in about
$0.5^2/\num{e-6} = \qty{2.5e5}{s}$, 2.9 days; the \qty{10}{m} \omterm{def:b2:magma-emplacement:bodies}{dyke} in $5^2/\num{e-6} =
\qty{2.5e7}{s}$, 289 days. The basalt (\qty{1.1}{h}) and the rhyolite in the wide
\omterm{def:b2:magma-emplacement:bodies}{dyke} (4.7 days) arrive before their \omterm{def:b2:magma-emplacement:bodies}{dykes} freeze; the rhyolite in the
\qty{1}{m} \omterm{def:b2:magma-emplacement:bodies}{dyke} (472 days) cannot.
\end{solution}

\begin{solution}{exo:b2:magma-emplacement:7}
Per kilogram of \omterm{def:b2:magma-emplacement:magma}{magma}, $n = 10/18.015 = \qty{0.555}{mol}$ of water and $nRT =
0.555\times8.314\times1173 = \qty{5414}{J}$; the remaining melt occupies $0.99/2400 =
\qty{4.13e-4}{m^3}$. At \qty{100}{MPa}: $V = \qty{5.41e-5}{m^3}$, fraction
$5.41/(5.41 + 41.3) = 0.12$. At \qty{10}{MPa}: $V = \qty{5.41e-4}{m^3}$, fraction
$0.57$. At \qty{0.1}{MPa}: $V = \qty{0.0541}{m^3}$, fraction $0.992$. The value at
\qty{100}{MPa} is the least reliable: water vapour at such pressure is dense and
far from ideal (and much of the water would in fact still be dissolved).
\end{solution}

\begin{solution}{exo:b2:magma-emplacement:8}
(a) $2300 + 60z = 2650$ gives $z = \qty{5.8}{km}$. (b) The melt has
\qty{2580}{kg/m^3}: $z = 280/60 = \qty{4.7}{km}$, about \qty{1.2}{km} higher. (c)
At $z = \qty{7.8}{km}$ the rock has $2300 + 60\times7.8 = \qty{2770}{kg/m^3}$; the
force is $(2770 - 2650)\times9.81 = \qty{1.2}{kN}$ per cubic metre, upward.
\end{solution}

\begin{solution}{exo:b2:magma-emplacement:9}
$v = 2\times300\times9.81\times(1000)^2/(9\times\num{e19}) = \qty{6.5e-11}{m/s}$, or
$\num{6.5e-11}\times\num{3.156e7} = \qty{2.1e-3}{m}$ a year, \qty{2.1}{mm/yr}. Rising
\qty{10}{km} takes $\num{e4}/\num{2.1e-3} = \qty{4.8}{Myr}$. The cooling time scale
is $(1000)^2/\num{e-6} = \qty{e12}{s} = \qty{32}{kyr}$: the body solidifies after
rising about $32000\times0.0021 \approx \qty{70}{m}$. A magmatic \omterm{def:b2:magma-emplacement:ascent}{diapir} cannot cross
such a crust; it needs hotter, softer surroundings.
\end{solution}

\begin{solution}{exo:b2:magma-emplacement:10}
(a) $x_{\mathrm{OH}} = \sqrt{K_2x_m} = \sqrt{K_1K_2f}$, so $x = K_1f + \frac12\sqrt{K_1K_2f}$.
The ratio of the first term to the second is $2\sqrt{K_1f/K_2}$: it tends to zero
when $f \to 0$ (the square root dominates) and grows without bound with $f$ (the
linear term dominates). (b) $K_1 f = \frac12\sqrt{K_1K_2 f}$ gives $K_1^2f^2 =
K_1K_2f/4$, so $f = K_2/(4K_1)$. (c) $f = 0.2/(\num{4e-4}) = \qty{500}{MPa}$. At
\qty{100}{MPa}: $x_m = 0.0100$, $\frac12\sqrt{\num{e-4}\times0.2\times100} =
\frac12\times0.0447 = 0.0224$, so $x = 0.0324$ and the hydroxyl share is
$0.0224/0.0324 = \qty{69}{\percent}$.
\end{solution}

\begin{solution}{exo:b2:magma-emplacement:11}
(a) $1000/\num{4e6} = \qty{2.5e-4}{km^3}$ a year, that is $\qty{2.5e5}{m^3}$ a year,
or $\num{2.5e5}/\num{3.156e7} = \qty{0.0079}{m^3/s}$, about eight litres a second.
(b) The \omterm{def:b2:magma-emplacement:bodies}{dyke} carries $1000\times2\times0.8 = \qty{1600}{m^3/s}$, and delivers
$\qty{e12}{m^3}$ in \qty{6.25e8}{s}, about 20 years: $\num{5e-6}$ of the
\qty{4}{Myr}. (c) Each sheet holds $20\times\num{e4}\times\num{e4} = \qty{2e9}{m^3} =
\qty{2}{km^3}$, so 500 sheets, one every $\num{4e6}/500 = 8000$ years on average.
(d) $t = (10)^2/(4\times0.65^2\times\num{0.7e-6}) = \qty{8.5e7}{s}$, 2.7 years. Each
sheet solidifies thousands of years before the next arrives: unless the supply
comes in rapid bursts, the \omterm{def:b2:magma-emplacement:bodies}{pluton} is never one large body of \omterm{def:b2:magma-emplacement:magma}{magma}, as the
dating of large \omterm{def:b2:magma-emplacement:bodies}{plutons} suggests.
\end{solution}

\begin{solution}{exo:b2:magma-emplacement:12}
(a) At fixed $x$, $\ln(T - T_0) = \text{const} - \frac12\ln t - x^2/(4\kappa t)$, whose
derivative $-1/(2t) + x^2/(4\kappa t^2)$ vanishes at $t^* = x^2/(2\kappa)$. Then
$e^{-x^2/(4\kappa t^*)} = e^{-1/2}$ and $\sqrt{4\pi\kappa t^*} = x\sqrt{2\pi}$, so $\Delta
T_{\max} = (T_m - T_0)\,b\,e^{-1/2}/(x\sqrt{2\pi}) = (T_m - T_0)\,b/(x\sqrt{2\pi e})$. (b)
$\sqrt{2\pi e} = 4.13$: $x = 1100\times2/(100\times4.13) = \qty{5.3}{m}$, reached after
$t^* = 5.3^2/(2\times\num{0.7e-6}) = \qty{2.0e7}{s}$, about 230 days (at only 2.7
\omterm{def:b2:magma-emplacement:bodies}{dyke} widths the plane-source approximation is rough). (c) The heating falls off
as $b/x$: a thin \omterm{def:b2:magma-emplacement:bodies}{dyke} carries little heat, and the rock is warmed appreciably
only within a few \omterm{def:b2:magma-emplacement:bodies}{dyke} widths.
\end{solution}

\begin{solution}{pb:b2:magma-emplacement:1}
\textbf{1.} $\log_{10}\eta = -4.55 + 5510/(1423 - 611) = -4.55 + 6.786 = 2.24$;
$\eta = \qty{172}{Pa.s}$.

\textbf{2.} $-4.55 + 5510/762 = -4.55 + 7.231 = 2.68$, $\eta = \qty{480}{Pa.s}$: a
factor of $480/172 = 2.8$ for \qty{50}{K}.

\textbf{3.} $172\times(1 - 1.35\times0.20)^{-2.5} = 172\times0.73^{-2.5} = 172\times2.20 =
\qty{378}{Pa.s}$.

\textbf{4.} $\phi = 1/1.35 = 0.74$. The field evidence places the end of bulk flow
at a melt fraction of \qtyrange{30}{50}{\percent}, that is
\qtyrange{50}{70}{\percent} of crystals: earlier, because real crystals are not
spheres of one size.

\textbf{5.} Basement: $2800 - 2710 = \qty{90}{kg/m^3}$, the \omterm{def:b2:magma-emplacement:magma}{magma} is lighter.
Sediments: $2450 - 2710 = \qty{-260}{kg/m^3}$, the \omterm{def:b2:magma-emplacement:magma}{magma} is denser.

\textbf{6.} $p = 2450\times9.81\times3000 = \qty{72.1}{MPa}$.

\textbf{7.} $90\times9.81 = \qty{883}{Pa/m}$.

\textbf{8.} $\bar u = 883\times2.0^2/(12\times378) = 3532/4536 = \qty{0.78}{m/s}$.

\textbf{9.} $q = 0.78\times2.0 = \qty{1.56}{m^2/s}$; over \qty{1.0}{km}, about
\qty{1560}{m^3/s}.

\textbf{10.} $Re = 2710\times0.78\times2.0/378 = 11$: laminar.

\textbf{11.} $12000/0.78 = \qty{15400}{s}$, about \qty{4.3}{h}.

\textbf{12.} $(1.0)^2/\num{0.7e-6} = \qty{1.43e6}{s}$, 16.5 days, far longer than
\qty{4.3}{h}: the \omterm{def:b2:magma-emplacement:bodies}{dyke} stays open while the \omterm{def:b2:magma-emplacement:magma}{magma} flows.

\textbf{13.} At \qty{3.0}{km} the \omterm{def:b2:magma-emplacement:magma}{magma} meets sediments lighter than itself: the
base of the sediments is its \omterm{def:b2:magma-emplacement:ascent}{level of neutral buoyancy}. Above it the \omterm{def:b2:magma-emplacement:magma}{magma} would
be denser than its surroundings, so it spreads sideways along the contact as a
\omterm{def:b2:magma-emplacement:bodies}{sill}.

\textbf{14.} $p_s = (0.40/0.34)^2 = \qty{1.38}{MPa}$, reached at
$\num{1.38e6}/(2450\times9.81) = \qty{58}{m}$ below the surface.

\textbf{15.} $0.34\times\sqrt{72.1} = 0.34\times8.49 = \qty{2.9}{\percent}$, far more than
the \qty{0.40}{\percent} present: the \omterm{def:b2:magma-emplacement:bodies}{sill} \omterm{def:b2:magma-emplacement:magma}{magma} exsolves no water.

\textbf{16.} CO$_2$ is much less soluble than water: it saturates and forms bubbles
deep in the ascent, so the first gas of a basalt is mostly CO$_2$.

\textbf{17.} Per kilogram, $n = 4.0/18.015 = \qty{0.222}{mol}$, $V = 0.222\times8.314\times1423/
\num{e5} = \qty{0.0263}{m^3}$; the melt occupies $0.996/2710 = \qty{3.68e-4}{m^3}$; the
fraction is $0.0263/(0.0263 + 0.000368) = 0.986$, \qty{98.6}{\percent}.

\textbf{18.} $(1150 + 100)/2 = \qty{625}{\celsius}$.

\textbf{19.} After 10 years, $2\sqrt{\kappa t} = 2\sqrt{\num{0.7e-6}\times\num{3.156e8}} =
\qty{29.7}{m}$. The point is at $x = \qty{60}{m}$ from the middle: $T = 100 +
525\,[\operatorname{erf}(-10/29.7) + \operatorname{erf}(110/29.7)] = 100 + 525\,(-0.365 +
1.000) = \qty{433}{\celsius}$.

\textbf{20.} At the middle, $T = 100 + 1050\operatorname{erf}(50/2\sqrt{\kappa t})$: after
10 years $1050\operatorname{erf}(1.68) = 1050\times0.983$, $T = \qty{1132}{\celsius}$; after
50 years, $2\sqrt{\kappa t} = \qty{66.5}{m}$, $1050\operatorname{erf}(0.752) = 1050\times0.713$,
$T = \qty{848}{\celsius}$.

\textbf{21.} At $x = \qty{130}{m}$: \qty{294}{\celsius} after 300 years,
\qty{295}{\celsius} after 363 years (the maximum), \qty{295}{\celsius} after 420
years. Without latent heat, the rock \qty{80}{m} above the \omterm{def:b2:magma-emplacement:bodies}{sill} never quite
reaches \qty{300}{\celsius}.

\textbf{22.} $\num{9e8}\times1.772/(\num{2.5e6}\times1050) = 0.608$.

\textbf{23.} $\lambda = 0.64$: $e^{-0.4096}/(0.64\times1.6346) = 0.6639/1.0461 = 0.635$;
$0.65$: $0.6554/1.0673 = 0.614$; $0.66$: $0.6469/1.0886 = 0.594$. The value $0.608$
lies between $0.65$ and $0.66$, close to $0.65$: $\lambda = 0.65$.

\textbf{24.} $100 + 1050/(1 + 0.642) = 100 + 640 = \qty{740}{\celsius}$. The latent
heat released at the solidification front keeps the \omterm{def:b2:magma-emplacement:magma}{magma} side at
\qty{1150}{\celsius} longer, so the contact receives more heat.

\textbf{25.} $(50)^2/\num{0.7e-6} = \qty{3.57e9}{s}$, 113 years.

\textbf{26.} $t = (100)^2/(16\times0.65^2\times\num{0.7e-6}) = \num{e4}/\num{4.73e-6} =
\qty{2.11e9}{s}$, so $\num{2.11e9}/\num{3.156e7}$: the \omterm{def:b2:magma-emplacement:bodies}{sill} is solid after
\textbf{\qty[text-series-to-math=true]{67}{} years}. It is shorter than
$(b/2)^2/\kappa$ by the factor $1/(4\lambda^2) = 0.59$: that time scale only sets
the order of magnitude, and both walls of the \omterm{def:b2:magma-emplacement:bodies}{sill} draw heat into a cold host at
once, fast enough to remove the latent heat as well.
\end{solution}
